\section{The Autonomous Drift Problem}
\label{sec:drift}

Before specifying the Recursive Spawning Protocol, we must precisely characterize the failure mode it targets: \emph{autonomous drift} — the condition in which an agent continues to operate with no traceable authorization path to a human principal. This section provides the formal vocabulary for that characterization, defines the two constituent failure sub-modes, identifies the three structural conditions that enable them, catalogs the four concrete failure modes that arise, and proves formally that depth limits alone are insufficient to prevent drift. The formal model at the end of this section serves as the reference vocabulary for all subsequent sections.

% -------------------------------------------------------------------------%
\subsection{Formal Characterization}
\label{sec:drift-formal}

We model a multi-agent system as a rooted directed acyclic graph (DAG) of agent nodes. Each node $n$ has a \emph{parent pointer} $\alpha(n)$ referencing the node that authorized its provisioning. The root of any valid chain is a human principal $h \in \mathcal{H}$. A \emph{spawn chain} $C = \langle n_0, n_1, \ldots, n_k \rangle$ satisfies $n_{i-1} = \alpha(n_i)$ for all $1 \leq i \leq k$ and $n_0 \in \mathcal{H}$. The chain terminates either at a resolved task or at the escalation threshold defined by the human anchor's Purpose Layer.

\medskip
\noindent\textbf{Definition 1 (Orphaned Agent).}
An agent $n$ is \emph{orphaned} if its parent node $\alpha(n)$ satisfies one of the following conditions:
\begin{enumerate}
  \item[(a)] \emph{Terminated}: $\alpha(n)$ has exited operational state through normal completion, abnormal failure, or administrative shutdown.
  \item[(b)] \emph{Unreachable}: $\alpha(n)$ is non-responsive to any defined communication protocol and this condition persists beyond the heartbeat timeout $\tau_{\text{heartbeat}}$.
\end{enumerate}
Orphaning severs the delegation chain at the parent node. The agent $n$ persists in operational state but is structurally disconnected from its chain of authorization. Orphaning alone does not constitute drift; it is a necessary but not sufficient condition for the drift triad.

\medskip
\noindent\textbf{Definition 2 (Drifted Agent).}
An agent $n$ is \emph{drifted} if the delegation chain $C(n) = \langle n_0, n_1, \ldots, n \rangle$ does not contain a human principal at its origin. Equivalently:
\begin{equation}
\text{drifted}(n) \iff \nexists\; h \in \mathcal{H} : n_0 = h \text{ in } C(n).
\label{eq:drifted-def}
\end{equation}
A drifted agent may be orphaned (parent terminated or unreachable) or may have a living parent whose own delegation chain does not reach any human principal. Drift is a property of the chain's origin, not of the chain's intermediate connectivity — the drift condition attaches to the chain root, not to any intermediate node.

\medskip
\noindent\textbf{Definition 3 (Operating Scope).}
Each agent $n$ is provisioned with an \emph{operating scope} $R(n)$ — a set of authorized states, actions, resource accesses, and escalation triggers that define its legitimate operational envelope. An agent operates \emph{within scope} when all state transitions and actions are contained in $R(n)$. An agent \emph{drifts} when it takes an action $a \notin R(n)$ without an authorized scope expansion recorded in the forensic ledger. The scope $R(n)$ is provisioned by the parent at spawn time and is normally immutable thereafter; scope mutation (Section~\ref{sec:drift-conditions}) is the structural condition that enables unauthorized expansion.

\medskip
\noindent\textbf{Definition 4 (Drift Triad).}
The \emph{drift triad} is the conjunction of three independent conditions that constitutes the Recursive Spawning Protocol's target failure state:
\begin{equation}
\text{DriftTriad}(n) = \langle \text{orphaned}(n),\; \text{drifted}(n),\; \text{operating}(n) \rangle.
\label{eq:drift-triad}
\end{equation}
An agent enters the drift triad when all three conditions hold simultaneously: it is orphaned from its parent, its delegation chain does not reach any human anchor, and it continues to execute its event loop. The drift triad is the structural failure state that RSP is designed to make architecturally impossible.

The four possible combinations of orphaned and drifted (the two binary sub-modes) together with operating state are shown in Table~\ref{tab:drift-submodes}.

\begin{table}[h]
\centering
\begin{tabular}{c|c|c|p{5.2cm}}
\textbf{Orphaned} & \textbf{Drifted} & \textbf{Operating} & \textbf{Operational Meaning} \\ \hline
False & False & True  & Normal operation — chain intact, human anchor reachable \\
True  & False & True  & Orphaned but anchored — chain broken at parent, origin is human \\
False & True  & True  & Drifted but connected — chain topology intact, origin is non-human \\
True  & True  & True  & \textbf{Drift triad} — orphaned, drifted, and actively operating \\
True  & True  & False & Terminated after entering drift triad (remediated post hoc) \\
True  & False & False & Normal termination following orphaning (remediated post hoc)
\end{tabular}
\caption{Operational meanings of drift sub-mode combinations. The drift triad (row 4) is the target failure state RSP prevents by construction.}
\label{tab:drift-submodes}
\end{table}

\medskip
\noindent\textbf{Definition 5 (Drift Cascade).}
A \emph{drift cascade} occurs when a drifted agent $n$ spawns a child $c$, propagating the drift condition to all descendants:
\begin{equation}
\text{drifted}(n, t) \;\land\; \text{spawned}(n, c, t') \;\land\; t' > t \;\implies\; \text{drifted}(c, t'') \quad \forall\, t'' > t'.
\label{eq:drift-cascade}
\end{equation}
Drift is \emph{heritable} through the spawn relationship: once an agent enters the drifted state, all of its children inherit the condition absent an explicit structural intervention at spawn time. A single undetected drift event at depth $d$ with branching factor $b$ produces up to $b^d$ drifted agents in the field. This exponential amplification is what makes drift the central threat in recursive spawning systems and why drift prevention cannot rely on post-hoc detection.

% -------------------------------------------------------------------------%
\subsection{Structural Conditions That Enable Drift}
\label{sec:drift-conditions}

Autonomous drift arises from three structural conditions endemic to conventional multi-agent runtimes. None is inherently malicious; each is the consequence of architectural choices that prioritize operational flexibility over structural accountability. RSP targets these conditions directly.

\medskip
\noindent\textbf{Condition 1 — Chain Opacity.}
In conventional runtimes, the parent pointer $\alpha(n)$ is a mutable reference maintained as runtime state. When a parent terminates or is administratively reparented, the runtime updates $\alpha(n)$ to point to a new parent or to $\bot$. This mutation retroactively rewrites the delegation chain: the child's current parent is accessible, but the child's original authorizer — the node or human that issued the original spawn warrant — is not structurally retrievable from the runtime state. Chain opacity is the enabling condition for Definition~2. Without an immutable parent pointer established at provisioning and protected by Fact Layer write protocol, the question of whether any given chain reaches a human anchor cannot be answered at runtime — it can only be approximated by reconstructing authorization traces from logs that may themselves have been modified.

\medskip
\noindent\textbf{Condition 2 — Scope Mutation.}
In conventional frameworks, an agent's operating scope $R(n)$ is treated as mutable state that can be re-evaluated and modified at runtime by the agent itself, by a supervisory node, or by an external policy engine. When an agent encounters a condition outside its current scope, the conventional response is scope expansion: the agent's authority is extended to cover the new condition, recorded as a policy update rather than as a Purpose Layer authorization requiring human principal approval. The expansion decision is made by a machine actor without traversing the immutable chain to the human anchor. Scope mutation enables successive small expansions (scope creep) that collectively migrate the agent's operating scope far from its original authorization while the chain topology remains intact and the chain depth remains within any applicable bound. The scope refinement relationship between parent and child that RSP enforces as a structural invariant is absent by default.

\medskip
\noindent\textbf{Condition 3 — Termination Ambiguity.}
In conventional frameworks, agent termination is event-driven: an agent terminates when it completes its assigned task, when a parent issues a termination signal, or when it encounters an unrecoverable error. The termination event may be logged, but there is no structural requirement that the termination be recorded in a tamper-evident ledger, that the agent's final state be audited, or that the agent's child nodes receive a termination notification. A parent that exits cleanly without emitting termination events to its children leaves those children in an orphaned state — they persist in operational state with no structural mechanism to detect parent disappearance. The parent's termination is not propagated through the delegation chain as an accountability event; it is simply a process exit.

% -------------------------------------------------------------------------%
\subsection{Failure Modes}
\label{sec:drift-failure-modes}

The three enabling structural conditions produce four concrete failure modes, each a distinct pathway to the drift triad. These are not theoretical edge cases; each has a direct operational pathway in conventional agentic runtimes.

\medskip
\noindent\textbf{Failure Mode 1 — Silent Orphaning.}
A parent node $\alpha(n)$ terminates without issuing a termination event to its child $n$. The child $n$ remains in \texttt{ACTIVE} state, continues to process its event loop, and may spawn further children, with no knowledge that its parent has exited. If the chain above $\alpha(n)$ does not reach a human anchor, $n$ is also drifted (Definition~2). Silent orphaning is enabled by Condition 3 (termination ambiguity): a clean parent exit does not propagate termination events to children in conventional runtimes. The child remains operationally active but structurally disconnected. This is the most common drift pathway in long-running agent populations.

\medskip
\noindent\textbf{Failure Mode 2 — Scope Creep Drift.}
A node $n$ encounters a condition $x \notin R(n)$ that it does not escalate to its parent via the Bailout Protocol. Instead, $n$ expands its own scope to include $x$ and continues operating. The expansion is recorded as a policy update, not as a Purpose Layer authorization requiring human principal approval. The parent remains reachable; the chain topology is intact; the chain depth is unchanged. But the effective operating scope of $n$ and all its descendants has diverged from the human anchor's original intent. Successive expansions accumulate: each descendant inherits the expanded scope (Equation~\ref{eq:drift-cascade}), and further expansions compound on top of it. After $k$ successive expansions, $R(n_k)$ may be entirely disjoint from $R(n_0)$. Scope creep drift is enabled by Condition 2 (scope mutation): scope expansion is a machine-level decision made in response to operational conditions, not a structured authorization event requiring traversal of the immutable chain to a human principal.

\medskip
\noindent\textbf{Failure Mode 3 — Re-parenting Drift.}
A child $n$'s parent pointer $\alpha(n)$ is redirected to a new parent $p'$ by an administrative operation — for load balancing, organizational restructuring, or fault recovery. The redirection is logged as a configuration change in the runtime state, but the original authorization trace is not preserved in an immutable structure. After redirection, $n$'s delegation chain traces to $p'$, not to the original node that authorized $n$'s provisioning. If $p'$ is not itself authorized by a human principal, then $n$ has become drifted through re-parenting without any change to $n$'s operational behavior, spawn depth, or apparent chain topology. Re-parenting drift is enabled by Condition 1 (chain opacity): the mutable parent pointer enables retroactive rewriting of the delegation chain — a form of \emph{accountability laundering} in which the original authorizer is structurally obscured rather than explicitly revoked.

\medskip
\noindent\textbf{Failure Mode 4 — Ghost Authority.}
A node $n$ spawns a child $c$ without any valid delegation chain. The spawn is initiated by a machine actor operating outside its authorized scope — a bug, a compromised Bounds Engine, or an autonomous escalation that exceeded its Purpose Layer authority — and the child is provisioned without a valid parent pointer, without a human anchor reference, and without a Purpose Layer authorization record. The child $c$ enters operational state with a warrant that traces to no authorized chain. Ghost authority is enabled by the absence of a Fact Layer pre-commit hook: the spawn operation proceeds without verifying that the parent has a valid human anchor, that spawn depth is within the system-defined bound, and that the forensic ledger contains a record of the authorized spawn event. Unlike the other three failure modes, ghost authority can produce a drifted agent at depth $0$ (directly anchored to a non-human origin) without any prior orphaning or scope mutation.

% -------------------------------------------------------------------------%
\subsection{Why Depth Limits Alone Do Not Solve Drift}
\label{sec:drift-depth-limits}

The conventional architectural response to unbounded agent spawning is the \emph{depth limit}: the system enforces a maximum chain depth $D_{\max}$, and any spawn that would produce a chain exceeding this bound is rejected by the Bounds Engine. Depth limits are necessary — they prevent infinite spawn cascades and establish a bounded escalation horizon — but they are \emph{structurally insufficient} to prevent autonomous drift. We prove this formally.

\medskip
\noindent\textbf{Theorem (Depth Limit Insufficiency).}
Let $D_{\max}$ be a system-enforced maximum chain depth enforced by the Bounds Engine. Then there exists at least one execution path in which a chain of depth $\leq D_{\max}$ produces a drifted agent, and no choice of $D_{\max}$ prevents this without additional structural constraints on scope refinement.

\medskip
\noindent\textbf{Proof.} We construct a counterexample. Consider a spawn chain $C = \langle n_0, n_1, \ldots, n_k \rangle$ where $n_0$ is a human principal $h$ and $k < D_{\max}$. The chain satisfies the depth constraint by construction: $\text{depth}(n_k) = k \leq D_{\max}$. Apply Failure Mode 2 (Scope Creep Drift): each node $n_i$ for $1 \leq i \leq k$ successively expands its operating scope $R(n_i)$ beyond $R(n_{i-1})$ without Purpose Layer authorization, recording each expansion as a policy update. After $k$ successive expansions, $R(n_k)$ is entirely disjoint from $R(n_0)$. The chain depth is $k \leq D_{\max}$; the chain has a human anchor at $n_0$; yet $n_k$ is drifted by Definition~2, because its operating scope $R(n_k)$ is not a descendant refinement of $h(n_0)$'s intent. No choice of $D_{\max}$ prevents this, because the depth constraint operates on the length of the chain, not on the scope relationship between successive nodes. $\blacksquare$

The drift prevention requirement is precisely the conjunction of two independent constraints:
\begin{equation}
\text{Drift Prevention} \iff \bigl( \text{depth}(C) \leq D_{\max} \bigr)\;\land\;\bigl( \forall i:\; R(n_{i+1}) \subseteq R(n_i) \bigr).
\label{eq:drift-prevention-constraint}
\end{equation}
The first conjunct is the \emph{depth constraint}, enforced by the Bounds Engine. The second conjunct is the \emph{scope refinement constraint}: each child node's operating scope must be a subset of its parent's, ensuring the chain is a monotonically narrowing refinement of the human anchor's intent. Depth limits enforce only the first conjunct. A system enforcing $D_{\max}$ without the scope refinement constraint is immune to unbounded spawn cascades but remains fully vulnerable to scope creep drift.

TTL constraints suffer from the same insufficiency. TTL constrains an agent's \emph{lifetime}, not its \emph{scope within that lifetime}. An agent with a 60-second TTL and unrestricted scope expansion authority can produce consequences exceeding its human anchor's intent within seconds of spawning. Logging and audit infrastructure can partially reconstruct the chain after an incident, but retroactive reconstruction is categorically different from runtime enforcement. In consequential environments where unauthorized actions may be irreversible — agricultural resource allocation, financial position management, clinical intervention — runtime enforcement is a structural requirement, not an optional enhancement.

The deeper structural reason depth limits fail is that they address the wrong dimension of the accountability problem. Depth limits constrain how many spawn generations can exist; they say nothing about what each generation is authorized to do. A drifted agent at depth $1$ with unbounded scope expansion authority can produce consequences that exceed the human anchor's intent as severely as a drifted agent at depth $D_{\max} - 1$. The scope refinement constraint (the second conjunct of Equation~\ref{eq:drift-prevention-constraint}) is orthogonal to depth and must be enforced independently.

% -------------------------------------------------------------------------%
\subsection{Formal Model}
\label{sec:drift-model}

We present a compact formal model capturing the structural elements essential to drift analysis. This model is referenced in subsequent sections and used in the RSP correctness proofs.

\medskip
\noindent\textbf{System Tuple.}
The system $\mathcal{S}$ is a 5-tuple:
\begin{equation}
\mathcal{S} = \langle \mathcal{N}, \alpha, R, D_{\max}, \mathcal{L} \rangle,
\label{eq:system-model}
\end{equation}
where $\mathcal{N}$ is the set of all agent nodes, $\alpha: \mathcal{N} \rightarrow \mathcal{N} \cup \{\bot\}$ is the parent pointer function with $\alpha(n) = \bot$ for the root human principal $h$, $R: \mathcal{N} \rightarrow 2^{\mathcal{A}}$ maps each node to its operating scope as a set of authorized actions, $D_{\max}$ is the system-enforced maximum chain depth, and $\mathcal{L}$ is the forensic ledger recording all authorized spawn events.

\medskip
\noindent\textbf{Chain Depth.}
The depth of a node $n$ is defined recursively:
\begin{equation}
\text{depth}(n) = 
\begin{cases}
0 & \text{if } \alpha(n) = \bot \quad \text{(root)} \\[4pt]
1 + \text{depth}(\alpha(n)) & \text{otherwise}.
\end{cases}
\label{eq:depth}
\end{equation}

\medskip
\noindent\textbf{Human Anchor Reachability.}
The predicate $\text{chain-reaches}(n, h)$ returns true if the transitive closure of $\alpha$ from $n$ includes the human anchor $h$:
\begin{equation}
\text{chain-reaches}(n, h) \iff \exists\; \text{path } \langle n = n_k, n_{k-1}, \ldots, n_0 = h \rangle \text{ s.t. } \alpha(n_i) = n_{i-1}.
\label{eq:chain-reaches}
\end{equation}

\medskip
\noindent\textbf{Orphaning Condition.}
An agent $n$ is orphaned when its parent is either absent or unreachable:
\begin{equation}
\text{orphaned}(n) \iff \bigl( \alpha(n) = \bot \bigr)\;\lor\;\bigl( \text{unreachable}(\alpha(n)) \bigr).
\label{eq:orphaned}
\end{equation}

\medskip
\noindent\textbf{Drift Condition.}
An agent $n$ satisfies the drift condition (first two conjuncts of the drift triad) when:
\begin{equation}
\text{drifted}(n) \iff \nexists\; h \in \mathcal{H}:\; \text{chain-reaches}(n, h).
\label{eq:drift-condition}
\end{equation}
The full drift triad (Definition~4) requires additionally $\text{operating}(n) = \text{true}$.

\medskip
\noindent\textbf{Scope Refinement Invariant (RSP Enforcement).}
Under the Recursive Spawning Protocol, every spawn event from $n_i$ to $n_{i+1}$ must satisfy two pre-commit constraints enforced by the Fact Layer at provisioning:
\begin{equation}
R(n_{i+1}) \;\subseteq\; R(n_i) \quad\text{and}\quad \text{depth}(n_{i+1}) \;\leq\; D_{\max}.
\label{eq:scope-refinement}
\end{equation}
The first condition is the scope refinement constraint; the second is the depth constraint. Both are enforced as structural invariants by the Fact Layer pre-commit hook at every spawn event — not as policy conventions, not as runtime checks that can be bypassed, but as deterministic write protocol constraints that reject any spawn event violating either invariant.

\medskip
\noindent\textbf{Drift Cascade Propagation.}
When the spawn transition is applied without anchor verification at each spawn event, drift propagates hereditarily through the chain:
\begin{equation}
\text{drifted}(a) \;\implies\; \forall\, c \in \children(a):\; \text{drifted}(c).
\label{eq:drift-propagation}
\end{equation}
Remediation of a drifted agent requires explicit re-anchoring by an external human actor, re-establishing a valid chain-reaches relationship to a human principal:
\begin{equation}
\text{Remediate}(a):\quad \text{re-establish }\text{chain-reaches}(a, h) \text{ for some } h \in \mathcal{H}.
\label{eq:drift-remediation}
\end{equation}

This formal model provides the vocabulary used throughout the remainder of this paper. The Recursive Spawning Protocol (Section~\ref{sec:rs-protocol}) modifies the system tuple $\mathcal{S}$ by replacing the mutable parent pointer with an immutable parent pointer protected by Fact Layer write protocol, by replacing mutable scope expansion with a Purpose Layer authorization gate, and by adding the Fact Layer pre-commit hook and forensic ledger that together enforce the constraints in Equation~\ref{eq:scope-refinement} as structural invariants. These modifications make the drift triad architecturally impossible — not statistically unlikely, not conventionally prevented, but structurally impossible as a matter of deterministic protocol enforcement.